\chapter{Discussion}
\label{chap:discussion}

The experimental results demonstrate that county-level wildfire Expected Annual Loss can be modeled with meaningful predictive power using accessible public data. A key finding is that environmental augmentation substantially improves performance over the county-only baseline. This suggests that loss-oriented wildfire modeling should not rely only on exposure, value, vulnerability, and socioeconomic indicators, but should also include hazard-relevant environmental conditions such as precipitation, vapor pressure deficit, vegetation composition, slope, elevation, and wind.

Unlike previous wildfire-risk studies that often evaluate predictive accuracy without translating results into consequence-oriented outputs, the present thesis converts model predictions into county-level wildfire Expected Annual Loss estimates, risk rankings, and exploratory ecological screening outputs.

The results support the central pipeline logic of the thesis. The FEMA National Risk Index provides a monetary wildfire-loss target, while the county-level predictors represent exposure, asset value, vulnerability, and resilience. Environmental augmentation then adds wildfire-hazard context derived from topography, land cover, and climate. Finally, machine-learning regression models and geographic validation convert these inputs into a county-level wildfire risk prediction and assessment workflow. In this way, the thesis moves beyond wildfire occurrence prediction and produces a consequence-oriented risk estimate in monetary units.

A central result is the improvement from the county-only baseline to the combined baseline-plus-environment model. The county-only model shows that exposure, value, vulnerability, and resilience variables contain useful predictive information. However, the much stronger performance of the environmental-only and combined models indicates that wildfire EAL is strongly shaped by environmental conditions. This supports the broader argument that wildfire losses should be understood as the result of interactions among hazard, exposure, vulnerability, and resilience rather than as a function of exposed assets alone.

The state-wise grouped validation provides an important perspective on geographic generalization. Although model performance decreases when entire states are held out, the combined model still retains substantial predictive power. This indicates that the model is learning a meaningful transferable wildfire-risk signal rather than merely benefiting from random geographic similarity between training and test observations. At the same time, the reduction in grouped-validation performance shows that regional generalization remains challenging. County-level wildfire-loss relationships vary across regions because of differences in fire regimes, vegetation structure, climate, settlement patterns, suppression systems, and local socioeconomic context.

The final risk-assessment output further strengthens the practical interpretation of the pipeline. By converting predicted log-transformed wildfire EAL back into monetary units, the model can rank counties by predicted wildfire loss and assign percentile-based risk categories. This step is important because it shows that the proposed workflow is not only a model-performance exercise. Instead, it produces an interpretable county-level risk-assessment product that can support comparison, prioritization, and further investigation of high-risk areas.

The exploratory ecological risk proxy adds a complementary screening layer to the monetary EAL model. By combining predicted wildfire-risk percentile with natural vegetation exposure, the analysis identifies counties where high predicted wildfire risk overlaps with high forest, shrub, and grass cover. This does not convert the thesis into a direct ecological-loss model, but it strengthens the ecological dimension by showing how the main FEMA/NRI pipeline can be extended toward ecological vulnerability screening. The economic--ecological overlap analysis further distinguishes counties with both high monetary risk and high ecological exposure from counties where only one of these dimensions is high.

Another important implication is conceptual. The experiments support the broader thesis argument that wildfire consequence assessment can be approached as a combined hazard--exposure--vulnerability problem. The county-only baseline captures part of the exposure and vulnerability signal, but the environmental augmentation shows that hazard-relevant climate, terrain, and land-cover information is essential for stronger loss-oriented prediction. This aligns with the literature gap identified in Chapter~\ref{chap:litreview}: wildfire-risk prediction becomes more useful when it is connected to expected consequences rather than treated as an endpoint by itself.

The thesis also contains an earlier Creek Fire tensor-based workflow as supporting work. This workflow was useful for developing a fine-scale data-engineering approach for wildfire modeling, including fire detections, fire perimeters, static geospatial predictors, and time-varying climate variables. However, it is not the main M.S. pipeline because it is not directly connected to monetary loss assessment. Therefore, it is placed in Appendix~\ref{app:creek_fire_pipeline} as a baseline approach that can support future fine-scale extensions of the FEMA/NRI loss-oriented pipeline.

The supporting Vietnam study provides an additional perspective on the broader wildfire-risk modeling direction of this thesis. While the FEMA/NRI pipeline focuses on county-level monetary Expected Annual Loss, the An Giang study focused on significant-fire classification using a compact deep neural network and multi-source environmental predictors. Its results showed strong classification performance, including 91.4\% accuracy, 74.8\% recall, 92.5\% precision, and 97.7\% specificity \citep{suliman2026angiang}. The study also showed that fire spread rate was the dominant SHAP feature, indicating the importance of dynamic fire-growth information for fine-scale wildfire prediction. Together, the Vietnam and FEMA/NRI results show two complementary directions: fine-scale wildfire event prediction and county-level monetary risk/loss assessment.

Despite the strengths of the FEMA/NRI pipeline, several limitations remain. First, the target variable is county-level Expected Annual Loss, which is spatially aggregated and does not represent event-level, parcel-level, or asset-level losses directly. Second, the environmental augmentation is restricted to the contiguous United States and excludes a small number of county-equivalent units because of geometry and dataset compatibility constraints. Third, the current modeling framework focuses on point prediction and feature importance; it does not yet include uncertainty quantification, calibrated prediction intervals, or probabilistic risk categories. Fourth, Random Forest feature importance provides useful first-order interpretation, but more detailed explainability methods such as SHAP could provide a deeper understanding of nonlinear feature interactions.

Although state-wise grouped validation is stricter than random splitting, it is not the only possible spatial validation strategy. Counties in different states can still share similar ecoregions, climate regimes, or vegetation types, while counties within large states may be highly heterogeneous. Future work could therefore evaluate transferability using ecoregion-based, climate-zone-based, or spatial-block validation strategies.

A further limitation concerns the ecological-loss dimension of the thesis. While ecological losses are discussed in the literature review and explored through the proxy-based screening layer, the empirical model does not directly predict ecological damage as a supervised target variable. The implemented pipeline focuses on monetary wildfire Expected Annual Loss from FEMA NRI. Environmental predictors such as forest fraction, shrub fraction, grass fraction, elevation, slope, precipitation, wind speed, and vapor pressure deficit provide ecological and hazard-related context, while the ecological proxy identifies overlap between predicted wildfire risk and natural vegetation exposure. However, this still does not constitute a direct ecological-loss model.

These limitations point toward several future extensions. Future work can improve geographic robustness, extend the environmental predictor set, add uncertainty quantification, and investigate richer representations of wildfire loss using event-level, asset-level, infrastructure, and ecological data. In particular, future ecological-loss modeling should incorporate explicit ecological vulnerability or recovery indicators, such as burn severity, ecosystem value, biodiversity exposure, protected-area overlap, vegetation recovery time, or post-fire resilience metrics. Nevertheless, the present results establish a defendable county-level wildfire loss-assessment pipeline. They show that publicly accessible exposure, vulnerability, resilience, topographic, land-cover, and climate datasets can be combined to produce meaningful wildfire-loss predictions in monetary units and exploratory ecological risk-screening outputs.
